\chapter{可视化与科研图表}

\section{学习目标}

学完本章后，你应该能够：

\begin{itemize}
    \item 理解科研图表的目标是传达结构、比较关系和支持论证，而不只是“把数据画出来”；
    \item 说清楚一张合格科学图至少应包含变量、单位、范围、标注和解读线索；
    \item 根据数据的数量级和关系特点，判断何时使用线性坐标、对数坐标、误差棒和标注；
    \item 从一个极小教学数据集生成可解释的散点图，并把图像与物理规律联系起来；
    \item 为后续的 HR 图、光谱图、时间域图和模型诊断图建立统一的图表思维方式。
\end{itemize}

进入本章前，建议你已经完成 Part 0 的基础训练，并能顺着读懂列表、循环、简单函数、文件读取和对象方法调用。本章会直接在这些能力之上讨论图像如何成为证据链的一部分。

\section{为什么科研图表不是“会用 matplotlib 就行”}

很多初学者会把绘图理解成最后一步：前面分析做完了，最后调用几条命令，把图画出来。但在科研里，图表并不是分析的装饰，而是分析论证的一部分。

一张图是否真正有用，取决于它能不能帮助读者回答以下问题：

\begin{itemize}
    \item 横轴和纵轴分别是什么物理量；
    \item 量纲和单位是否清楚；
    \item 图中的趋势究竟说明了什么关系；
    \item 数据散布是真实结构、观测误差，还是样本选择效应；
    \item 这张图是否足以支撑文本中的论断。
\end{itemize}

\begin{defbox}[科学图表]
科学图表是把数据、关系和解释组织成视觉证据的工具。它的目标不是单纯美观，而是让读者更高效地理解并检验一个科学论点。
\end{defbox}

因此，本章真正讨论的不是“如何调用作图库”，而是“怎样把一张图变成一条论证链的一部分”。

\section{从数据到图像：绘图前你其实已经做了很多决定}

在任何一张科研图出现之前，你至少已经做了四个层面的决定：

\begin{enumerate}
    \item \textbf{选择了哪些数据}：不是所有字段都应该进图；
    \item \textbf{选择了什么关系}：是比较 $x$ 与 $y$，还是比较残差与误差；
    \item \textbf{选择了什么尺度}：线性、对数，还是归一化；
    \item \textbf{选择了什么编码}：颜色、符号、误差棒、标注。
\end{enumerate}

换句话说，图并不是“中性结果展示”，而是研究者对数据关系的一次有意识表达。

\section{一个最小示例：为什么行星轨道图有教学价值}

本章继续使用一个小型行星教学数据表，包含半长轴和轨道周期。这个例子之所以适合 Part I，并不是因为它复杂，而正是因为它足够简单，让你能集中思考图表逻辑本身。

\begin{examplebox}[读取教学数据并准备变量]
\begin{lstlisting}[style=pythonstyle]
import csv
from pathlib import Path

data_path = Path("../../data/small/planetary_orbits_demo.csv")
with data_path.open(newline="", encoding="utf-8") as f:
    rows = list(csv.DictReader(f))

planet_names = [row["planet"] for row in rows]
x_values = [float(row["semi_major_axis_au"]) for row in rows]
y_values = [float(row["orbital_period_years"]) for row in rows]
\end{lstlisting}
\end{examplebox}

到这里为止，我们还没有开始画图，但已经完成了一个最小可视化算法的前三步：

\begin{itemize}
    \item 选择数据列；
    \item 明确横纵变量；
    \item 把文本数据转换为可计算的数值。
\end{itemize}

\section{变量、单位与标注：图表的最小语言}

一张科学图至少需要解决三种语言问题：

\subsection{变量语言}

读者必须知道横轴和纵轴代表什么量。例如本章中是：

\[
x = a \quad (\text{semi-major axis}),
\qquad
y = P \quad (\text{orbital period}).
\]

\subsection{单位语言}

变量名如果没有单位，解释会立刻变得含糊。例如：

\[
a \; [\mathrm{au}], \qquad P \; [\mathrm{yr}]
\]

和单写“semi-major axis”“orbital period”相比，前者更完整，因为它告诉读者数值的量纲。

\subsection{对象语言}

当样本很少且个体身份重要时，单靠点本身还不够，你往往需要标出对象名称。对于行星数据，这样做特别合理，因为每一个点都对应一个熟悉的天体，而不是一个大样本中的匿名成员。

\section{一张最小散点图是怎样构成的}

在 notebook 中，最小散点图代码如下：

\begin{examplebox}[绘制轨道周期与半长轴]
\begin{lstlisting}[style=pythonstyle]
fig, ax = plt.subplots(figsize=(6, 4))
ax.scatter(x_values, y_values, color="tab:blue")

for name, x_value, y_value in zip(planet_names, x_values, y_values):
    ax.annotate(name, (x_value, y_value), fontsize=8,
                xytext=(4, 4), textcoords="offset points")

ax.set_title("Orbital Period vs. Semi-major Axis")
ax.set_xlabel("Semi-major axis [au]")
ax.set_ylabel("Orbital period [yr]")
ax.grid(alpha=0.3)
\end{lstlisting}
\end{examplebox}

代码表面上只是几个函数调用，但每一行都对应一个图表设计决策：

\begin{itemize}
    \item \texttt{scatter} 用于强调离散样本；
    \item \texttt{annotate} 用于识别个体；
    \item 标题说明“比较关系是什么”；
    \item 轴标签说明“量是什么、单位是什么”；
    \item 网格线提高读图效率，但强度被压低，避免喧宾夺主。
\end{itemize}

\section{图表类型选择：关系决定图形}

初学者很容易把“画图”理解成“先试一个散点图”。散点图当然重要，但科研图表的第一步应该是判断你想表达什么关系。常见选择可以先记住下面这组最小对应：

\begin{center}
\begin{tabular}{lll}
\toprule
目标 & 常见图表 & 适合回答的问题 \\
\midrule
两个连续变量关系 & 散点图 & 是否有关联或趋势 \\
随时间或波长变化 & 折线图 & 结构是否沿有序轴变化 \\
样本分布 & 直方图 & 数值集中在哪些范围 \\
模型诊断 & 残差图 & 模型是否留下系统结构 \\
测量不确定度 & 误差棒图 & 点的位置是否足够可信 \\
\bottomrule
\end{tabular}
\end{center}

这张表不是固定规则，而是提醒你：图表类型应当服务于问题。HR 图强调颜色和绝对星等关系，适合散点图；光谱有自然波长轴，通常用折线图；周期搜索后的残差，则应该用诊断图来检查模型遗漏了什么。

\begin{protip}[先说问题，再选图]
如果你说不清一张图想回答什么问题，通常也很难选对图表类型。先写一句“这张图用来比较什么”，再写绘图代码。
\end{protip}

\section{为什么有时要使用对数坐标}

很多自然科学关系不是“加法式变化”，而是“乘法式变化”或幂律关系。此时，对数坐标非常重要。

例如，开普勒第三定律可以写成：

\[
P^2 \propto a^3.
\]

若取对数，就得到：

\[
\log P = \frac{3}{2} \log a + C,
\]

其中 $C$ 是常数。这说明：\textbf{如果原始关系是幂律，那么在对数坐标下会更接近直线。}

这也是为什么天文和物理中很多图都喜欢用对数轴。它不是为了显得高级，而是因为：

\begin{itemize}
    \item 大动态范围数据更容易看；
    \item 乘法关系更容易被识别；
    \item 不同数量级的对象可以放到同一张图里比较。
\end{itemize}

\section{图表不只要“画出来”，还要“读出来”}

对于本章这个最小行星数据集，图中最重要的信息并不是“点很多”，而是：

\begin{itemize}
    \item 半长轴越大，轨道周期越长；
    \item 这种增长不是简单线性，而更接近幂律关系；
    \item 小样本图像仍然可以服务于物理判断，而不只是数据展示。
\end{itemize}

这也是教学中的关键一跃：画图从“输出结果”变成“组织直觉”。后面不论是 HR 图、回归残差图、ROC 曲线还是光谱线型图，本质上都要回到这一点。

\section{从图到论断：claim--evidence--limit}

科研图表真正进入报告或论文时，不能只停留在“如图所示”。一张图应当至少支持三句话：

\begin{enumerate}
    \item \textbf{论断（claim）}：你希望读者从图中相信什么；
    \item \textbf{证据（evidence）}：图中的哪些结构支持这个判断；
    \item \textbf{限制（limit）}：这个判断在哪些条件下成立，哪里还不能过度解释。
\end{enumerate}

对行星轨道示例来说，一个合格的读图表述可以是：

\begin{examplebox}[把图转成科学表述]
这张图支持“轨道周期随半长轴增大而增大，并且更接近幂律而不是线性关系”的判断。证据是内外行星在对数坐标下呈现近似直线趋势，且 $P^2/a^3$ 的数值在教学数据中接近常数。限制是样本只包含太阳系行星，不能直接推广到所有系外行星系统。
\end{examplebox}

这个三句式看起来朴素，但很适合训练初学者。第一句防止图变成装饰；第二句要求你真的回到数据结构；第三句则防止从一张小图跳到过大的结论。后面的 HR 图、光谱线型图、模型残差图和混淆矩阵，都可以用同样方式阅读。

\section{图表中的误差、比较与诊断}

一张真正有分析价值的科学图，往往不仅仅展示一个变量关系，还应该考虑：

\begin{itemize}
    \item 是否需要误差棒；
    \item 是否需要对照组；
    \item 是否需要拟合线、残差图或分组颜色；
    \item 是否需要说明哪些点是异常值、边界样本或低质量样本。
\end{itemize}

虽然本章还不要求一次覆盖所有图表类型，但必须让学生知道：\textbf{“散点图 + 标题”只是最小起点，不是全部。}

\section{误差棒与残差图：让图表承担诊断功能}

如果每个观测点都有不确定度，图中最好不要只画点。误差棒能直接告诉读者：一个点的位置大致可信到什么程度。

\begin{examplebox}[最小误差棒代码]
\begin{lstlisting}[style=pythonstyle]
ax.errorbar(
    x_values,
    y_values,
    yerr=y_errors,
    fmt="o",
    capsize=3,
    color="tab:blue",
)
\end{lstlisting}
\end{examplebox}

误差棒的意义不是让图“更专业”，而是让读者能判断趋势是否明显超过测量噪声。如果两个点的误差范围大幅重叠，那么单靠点的位置差异就不应该被过度解释。

\subsection{残差图为什么经常比主图更诚实}

当你有一个模型预测 $y_{\mathrm{model}}(x)$ 时，残差定义为：

\[
r_i = y_i - y_{\mathrm{model}}(x_i).
\]

主图常常让人先看到“模型大体贴合数据”，但残差图会更直接地暴露模型哪里没解释好。一个好的残差图通常应该检查：

\begin{itemize}
    \item 残差是否围绕零上下波动；
    \item 残差是否随 $x$ 出现系统趋势；
    \item 是否有少数异常点主导了拟合结果；
    \item 残差大小是否与误差棒量级一致。
\end{itemize}

\begin{examplebox}[最小残差计算]
\begin{lstlisting}[style=pythonstyle]
residuals = [
    y_obs - y_model
    for y_obs, y_model in zip(y_values, model_values)
]
\end{lstlisting}
\end{examplebox}

这也是后面机器学习和物理拟合章节会反复出现的习惯：不要只画预测值，也要画模型没有解释掉的部分。

\section{视觉编码也会影响科学判断}

图表的颜色、点大小、线宽和坐标范围并不是中性的装饰。它们会影响读者先看到什么、忽略什么，以及对差异大小的直觉判断。

例如，同一组数据如果把纵轴范围截得很窄，微小变化会显得非常剧烈；如果颜色映射没有说明单位和范围，读者可能会把颜色差异误读成物理分类；如果图例缺失，分组信息就很难被复查。

\begin{protip}[把视觉编码翻译回科学含义]
每使用一种颜色、符号、线型或点大小，都应该能回答一句话：它编码的是什么科学信息？如果回答不出来，这个视觉元素往往就是噪声。
\end{protip}

一个安全的自问流程是：

\begin{enumerate}
    \item 这个视觉元素对应哪个变量或分类；
    \item 它是否有图例、色标或文字说明；
    \item 它会不会夸大某个差异；
    \item 色盲读者或黑白打印时是否仍能理解主要结构；
    \item 图中最醒目的部分是否真的对应最重要的科学信息。
\end{enumerate}

这类问题在后面的 HR 图、光谱分类、异常检测和模型诊断里尤其重要。图表越漂亮，越要警惕它是否把读者带向了错误的解释。

\section{从绘图命令到最小工作流算法}

把本章内容压缩成一个最小算法，可以写成：

\begin{enumerate}
    \item 明确想比较的物理量；
    \item 检查它们的单位与数量级；
    \item 判断线性还是对数坐标更适合；
    \item 选择点、线、颜色、误差棒等图形编码；
    \item 添加标题、轴标签和必要标注；
    \item 保存图像并记录生成代码；
    \item 回到物理问题，解释图中结构意味着什么。
\end{enumerate}

真正“会画科研图”的标志，不是记住了多少 Matplotlib 参数，而是能把这条流程稳定走通。

\section{保存图像：输出文件也是科研记录}

一张图如果只存在于 notebook 当前输出里，就很难进入报告、论文、复现流程或版本控制。科研图表通常应该显式保存，并且文件名要能表达图的内容和版本。

\begin{examplebox}[保存图像的最小写法]
\begin{lstlisting}[style=pythonstyle]
fig.tight_layout()
fig.savefig("results/orbit_period_vs_axis.png", dpi=200)
fig.savefig("results/orbit_period_vs_axis.pdf")
\end{lstlisting}
\end{examplebox}

常见格式可以这样理解：

\begin{itemize}
    \item \texttt{png}：适合网页、课件和快速预览；
    \item \texttt{pdf}：适合论文、讲义和矢量输出；
    \item \texttt{svg}：适合可编辑矢量图，但要注意字体兼容；
    \item \texttt{jpg}：适合照片，不适合线图和文字较多的科研图。
\end{itemize}

\begin{warnbox}[图像文件名也要可复现]
不要把最终图命名为 \texttt{plot.png}、\texttt{new.png} 或 \texttt{final2.png}。更好的文件名应包含对象、变量关系、版本或日期线索，例如 \texttt{hr\_diagram\_quality\_cut\_v1.pdf}。
\end{warnbox}

\section{图表交付前的自审清单}

科研图表在进入报告、作业或论文草稿前，最好先做一次自审。一个最小清单可以包括：

\begin{enumerate}
    \item 轴标签是否包含变量名和单位；
    \item 图例是否解释了颜色、符号或线型；
    \item 坐标尺度是否适合数据动态范围；
    \item 是否需要误差棒、残差图或样本量说明；
    \item 文件名是否能追踪图的内容和版本；
    \item 图注是否说明了数据来源、筛选条件和主要结论。
\end{enumerate}

\begin{protip}[图注不是装饰]
图注应该帮助读者理解这张图支持什么判断。一个好图注通常包含：数据来源、关键处理、坐标含义、主要结构和限制条件。只写“某某关系图”通常不够。
\end{protip}

\section{一个最小图表说明模板}

图表进入报告前，最好配一段短说明，而不只是保存成文件。一个最小模板可以写成：

\begin{quote}
图 \texttt{orbit\_period\_vs\_axis.png} 使用 \texttt{planetary\_orbits\_demo.csv} 中的半长轴和轨道周期，横轴单位为 AU，纵轴单位为 year。对数坐标下点列近似满足 \(P^2/a^3\) 为常数的结构，支持 Kepler 第三定律的教学判断；该图只包含少量教学行星样本，不应用来讨论真实观测误差、样本选择或多体摄动。
\end{quote}

这段话把图表文件、数据来源、坐标单位、主要 claim 和限制条件放在一起。到了最终项目包，图表本身不是证据的全部；图表说明才把视觉结构连接到 Interpretation card，并把数据来源和生成入口连接到 Data card 与 Reproducibility card。

这个清单和模板会在后面的 HR 图、光谱、时间序列和机器学习诊断图中反复出现。也就是说，本章学到的不是 Matplotlib 技巧，而是科研证据表达的底层规范。

\section{配套 notebook 做了什么}

配套 notebook 延续 \texttt{planetary\_orbits\_demo.csv} 这个小数据集，从读取半长轴和轨道周期开始，逐步完成：

\begin{itemize}
    \item 检查数据点数量、变量范围和单位；
    \item 计算 $P^2/a^3$，把图像与 Kepler 关系联系起来；
    \item 比较线性直觉和对数坐标下的幂律结构；
    \item 构造一个最小模型并计算残差；
    \item 在有 \texttt{matplotlib} 的环境中生成图像，在没有图形库时仍保留数值解释。
\end{itemize}

这个设计故意不依赖复杂数据。它要训练的是一件更基础的事：一张图不能只是 notebook 输出，它必须能回答“变量是什么、单位是什么、关系是什么、模型哪里没有解释好”。

\section{常见错误}

\begin{warnbox}[科研绘图中的典型问题]
\begin{itemize}
    \item 没有标题或没有轴标签；
    \item 有变量名但没有单位；
    \item 量纲跨度很大却坚持用不合适的线性坐标；
    \item 一张图里元素太多，却没有图例、说明或视觉层次；
    \item 只把图当成装饰，不去解释图中的关系和结构。
\end{itemize}
\end{warnbox}

\section{AI 助手如何帮助你}

在这一章，AI 助手适合帮助你：

\begin{itemize}
    \item 帮你把基础绘图代码整理得更清楚；
    \item 提醒你是否遗漏了标题、单位、图例或误差说明；
    \item 根据数据动态范围建议是否尝试对数坐标；
    \item 帮你把图像描述改写成更有物理含义的文字说明。
\end{itemize}

但图表到底应该支持什么科学论点、某种坐标变换是否会误导读者、以及异常点应如何解释，仍然必须由你自己判断。

\section{练习}

\begin{exercisebox}[基础练习]
把纵轴改为对数坐标，并结合公式
\[
\log P = \frac{3}{2}\log a + C
\]
说明为什么对数图在幂律关系中更有解释力。
\end{exercisebox}

\begin{exercisebox}[进阶练习]
尝试将内行星和外行星用不同颜色显示，并写一段说明：这种分组编码到底在帮助读者看到什么，而不仅仅是在“区分颜色”。
\end{exercisebox}

\begin{exercisebox}[开放练习]
为你自己的一个科学小项目设计“绘图前检查清单”，至少包含：变量、单位、坐标尺度、图形编码、解释目标和可能误导读者的风险。
\end{exercisebox}

\section{本章小结}

科研图表不是把数据变成图，而是把数据变成可讨论、可比较、可检验的视觉证据。本章用最小行星例子说明了一个重要思想：图表的真正价值，在于帮助你和读者一起看见结构、组织直觉，并把后续分析建立在更清楚的证据表达上。
